
%(BEGIN_QUESTION)
% Copyright 2006, Tony R. Kuphaldt, released under the Creative Commons Attribution License (v 1.0)
% This means you may do almost anything with this work of mine, so long as you give me proper credit

Ethvert stoff som øker elektrisk ledningsevne når det løses i vann kalles en {\it elektrolytt}. Når bordsalt blandes i vann skjer noe interessant kalt {\it dissosiasjon}. En lignende prosess kalt {\it ionisering} skjer når hydrogenklorid (HCl) går i løsning. Forklar disse prosessene.

Identifiser deretter hvordan følgende elektrolytter deles i kationer og anioner i vann. Altså: hvilken del av hvert molekyl blir anion, og hvilken del blir kation:

\begin{itemize}
\item{} Natriumklorid (NaCl) $\to$
\vskip 5pt
\item{} Hydrogenklorid (HCl) $\to$
\vskip 5pt
\item{} Natriumhydroksid (NaOH) $\to$
\vskip 5pt
\item{} Kalsiumsulfat (CaSO$_{4}$) $\to$
\vskip 5pt
\item{} Kaliumsulfat (K$_{2}$SO$_{4}$) $\to$
\vskip 5pt
\item{} Svovelsyre (H$_{2}$SO$_{4}$) $\to$
\end{itemize}

\vskip 10pt

Hint: et ``Periodic Table of the Ions'' kan være nyttig for å avgjøre hvordan disse forbindelsene dissosierer.

\underbar{file i00604}
%(END_QUESTION)





%(BEGIN_ANSWER)

Både {\it dissosiasjon} og {\it ionisering} beskriver separasjon av tidligere bundne atomer ved oppløsning. Forskjellen ligger i stofftypen: ``dissosiasjon'' brukes om ioniske forbindelser (f.eks. bordsalt), mens ``ionisering'' brukes om kovalente (molekylære) forbindelser som HCl, som ikke er ioniske i ren tilstand.

\begin{itemize}
\item{} Natriumklorid (NaCl) $\to$ Na$^{+}$ + Cl$^{-}$
\vskip 5pt
\item{} Hydrogenklorid (HCl) $\to$ H$^{+}$ + Cl$^{-}$
\vskip 5pt
\item{} Natriumhydroksid (NaOH) $\to$ Na$^{+}$ + OH$^{-}$
\vskip 5pt
\item{} Kalsiumsulfat (CaSO$_{4}$) $\to$ Ca$^{2+}$ + SO$_{4}^{2-}$
\vskip 5pt
\item{} Kaliumsulfat (K$_{2}$SO$_{4}$) $\to$ 2K$^{+}$ + SO$_{4}^{2-}$
\vskip 5pt
\item{} Svovelsyre (H$_{2}$SO$_{4}$) $\to$ 2H$^{+}$ + SO$_{4}^{2-}$
\end{itemize}

\vskip 10pt

Kjennetegnet på en ionisk forbindelse er at den leder strøm i ren væskefase. Selv i fast form er en ionisk forbindelse allerede ionisert, med atomer bundet av ladningsubalanse. Væskefasen gir mobilitet nok til dissosiasjon.

\vskip 10pt

Molekylære (kovalente) forbindelser skiller seg normalt ikke lett i ioner i ren væskefase. Kovalente bindinger bygger på {\it delte} elektroner, ikke elektrostatisk tiltrekning mellom ioner. Et elektron må flyttes fra ett atom til et annet for at en kovalent forbindelse skal brytes i ioniserte deler, derfor kalles prosessen ionisering. Vann er et godt eksempel. Noe ionisering skjer, men andelen er svært liten, så rent vann regnes som elektrisk isolator.

Andre kovalente forbindelser finnes også, som hydrogenklorid. HCl ioniseres til hydrogenkationer og kloranioner i vann, og blir dermed en elektrolytt selv om ren HCl-væske ikke leder strøm.

\vskip 10pt

Ved elektrolyse av saltvann frigjøres både oksygen- og klorgass ved anoden (+), mens hydrogengass frigjøres ved katoden ($-$). Natrium samles ved katoden. Forholdet klor/oksygen ved anoden avhenger av saltkonsentrasjonen. Ved høy nok strøm kan kobberioner løses i vannet og gi blågrønn farge.

%(END_ANSWER)





%(BEGIN_NOTES)


%INDEX% Kjemi, grunnleggende: ioniske versus kovalente bindinger
%INDEX% Kjemi, ion: dissosiasjon

%(END_NOTES)

